% ID: SEL-D0018-BD-E07
% ESTADO: cotejada. Segunda pasada visual completa y borrador compilado revisados.
% Fuente: Selectividad/matematicas_aplicadas/matematicas_aplicadas_P_F.pdf
% SHA-256: fe6fd5b48407569d41169c80d71e6997050d6d06d1afd41cd639a553f9b712b5
% Página: 2; bloque D; ejercicio original 7.

Tomada al azar una muestra de 600 alumnos de una universidad española, se encontró que $2/3$ de los mismos podían expresarse en inglés con fluidez.
\begin{enumerate}[label=\alph*)]
\item \textbf{(1.5 puntos)} Calcule un intervalo de confianza al 98\% para estimar la proporción de alumnos de esa universidad que pueden expresarse en inglés con fluidez. \textquestiondown{}Se podría admitir a ese nivel de confianza que la proporción de alumnos de esa universidad que pueden expresarse en inglés con fluidez es $13/20$?
\item \textbf{(0.25 puntos)} Teniendo en cuenta el intervalo anterior, \textquestiondown{}qué error máximo se cometería en dicha estimación?
\item \textbf{(0.75 puntos)} Si se mantienen la misma proporción muestral y la misma confianza, \textquestiondown{}cuántos alumnos como mínimo habría de tener una muestra para que el error de estimación sea inferior al 2\%?
\end{enumerate}
